\section*{Chapter \ref{ch:g11:perfect-competition} --- Perfect Competition}
% figdata: figdata/grade-11/03-perfect-competition/out/entry.csv; Brenland numbers invented.

\begin{solution}{exo:g11:perfect-competition:1}
$100 \times 5 - 420 = 80$ dalens.
\end{solution}

\begin{solution}{exo:g11:perfect-competition:2}
$C'(q) = 2q = 10$ gives $q = 5$. Profit $50 - (20 + 25) = 5$ dalens.
\end{solution}

\begin{solution}{exo:g11:perfect-competition:3}
No: the price is below the lowest \omterm{def:g11:production-and-costs:atc}{average variable cost}, so every unit sold loses money on top of the fixed cost. Stopping, it loses its fixed cost, 200 dalens;
producing, it would lose more.
\end{solution}

\begin{solution}{exo:g11:perfect-competition:4}
Highest birth rate: computer programming, 15.2 per 100; lowest: legal and accounting, 5.8. Retail trade and computer programming, 8.0 each, had 8 or more closures per
100.
\end{solution}

\begin{solution}{exo:g11:perfect-competition:5}
\omterm{def:g11:perfect-competition:economic-profit}{Accounting profit} $80000 - 50000 = \dal{30000}$; \omterm{def:g11:perfect-competition:economic-profit}{economic profit} $30000 - 35000 = -\dal{5000}$, an economic loss.
\end{solution}

\begin{solution}{exo:g11:perfect-competition:6}
$\mathit{MC} = 2 + 2q$, $\mathit{AVC} = 2 + q$, lowest (2) at the smallest outputs: the \omterm{def:g11:perfect-competition:shutdown-price}{shutdown price} is 2. Supply: $2 + 2q = p$, so $q = (p - 2)/2$ for $p > 2$, and 0
below.
\end{solution}

\begin{solution}{exo:g11:perfect-competition:7}
$\mathit{AVC} = 4 - q + q^2/3$ has derivative $-1 + 2q/3$, zero at $q = 1.5$, where $\mathit{AVC} = 4 - 1.5 + 0.75 = 3.25$ (and $\mathit{MC}(1.5) = 4 - 3 + 2.25 = 3.25$).
Supply: $p = 4 - 2q + q^2$ gives $(q - 1)^2 = p - 3$, and on the rising part of \omterm{def:g11:marginal-reasoning:marginal-benefit}{marginal cost} $q = 1 + \sqrt{p - 3}$, for $p \geq 3.25$; 0 below.
\end{solution}

\begin{solution}{exo:g11:perfect-competition:8}
Industry supply $Q = 100(p - 2)/2 = 50(p - 2)$. $50p - 100 = 600 - 50p$ gives $p = 7$, $Q = 250$, and 2.5 units per firm.
\end{solution}

\begin{solution}{exo:g11:perfect-competition:9}
Long-run price 20: $Q = 10000 - 4000 = 6000$, so $6000/50 = 120$ firms. After the rise, $Q = 12000 - 4000 = 8000$, so 160 firms.
\end{solution}

\begin{solution}{exo:g11:perfect-competition:10}
$\mathit{ATC} = 18/q + 2q$, with derivative $-18/q^2 + 2$, zero at $q = 3$, where $\mathit{ATC} = 6 + 6 = 12$. Long-run price 12, each firm producing 3.
\end{solution}

\begin{solution}{exo:g11:perfect-competition:11}
Before: price 3, output 8. After: price 3.5, output 10. Profit $3.5 \times 10 - (8 + 10 + 12.5) = 4.5$ hundred dalens, 450 dalens a week; the \omterm{def:g11:production-and-costs:atc}{average total cost} at 10 is
3.05, and the shaded rectangle is $(3.5 - 3.05) \times 10 = 4.5$.
\end{solution}

\begin{solution}{exo:g11:perfect-competition:12}
High \omterm{def:g11:perfect-competition:economic-profit}{economic profits} lasting for many years in an industry where entry is free, with no sign of entry; or prices staying well above the lowest average cost of firms
that could enter. \omterm{def:g11:perfect-competition:economic-profit}{Accounting profits} leave out the owners' time and the return their money would earn elsewhere: a positive \omterm{def:g11:perfect-competition:economic-profit}{accounting profit} is what zero \omterm{def:g11:perfect-competition:economic-profit}{economic
profit} looks like in the accounts. A test needs a measure of those opportunity costs, for example the return on capital compared with the return available elsewhere.
\end{solution}

\begin{solution}{exo:g11:perfect-competition:13}
Entry raises the price of the shared input, which raises every firm's costs and the lowest \omterm{def:g11:production-and-costs:atc}{average total cost}. After a rise in demand, the industry settles with more
firms but at a higher lowest average cost, so the long-run price is higher. The long-run \omterm{def:g11:perfect-competition:industry-supply}{industry supply curve} slopes upwards instead of being flat at $m$.
\end{solution}

\begin{solution}{exo:g11:perfect-competition:14}
The number of taxis cannot rise, so fares rise with demand and each taxi earns more. The extra earnings cannot be competed away by entry; they go to whoever owns a
licence, and the licence price rises to the value of those earnings. The gain goes to the licence holders at the time of the rise; later buyers pay for it in the
licence price.
\end{solution}

\begin{solution}{exo:g11:perfect-competition:15}
A newcomer enters only if the profits it expects cover the entry cost, so the price can stay above break-even by up to that margin without attracting entry. A firm
already in has paid its entry cost, which is sunk: it stays as long as the price covers its \omterm{def:g11:production-and-costs:atc}{average variable cost}, even below break-even. Between those two prices,
neither entry nor exit happens.
\end{solution}

\begin{solution}{pb:g11:perfect-competition:1}
Quantities in hundreds of dozens a week; costs and profits in hundreds of dalens a week.

\textbf{1.} $\mathit{MC} = 1 + 0.25q$; $\mathit{AVC} = 1 + 0.125q$; $\mathit{ATC} = 8/q + 1 + 0.125q$.

\textbf{2.} $-8/q^2 + 0.125 = 0$ gives $q = 8$; $\mathit{ATC} = 1 + 1 + 1 = 3$.

\textbf{3.} \omterm{def:g11:production-and-costs:atc}{Average variable cost} is lowest, at 1, for the smallest outputs: the \omterm{def:g11:perfect-competition:shutdown-price}{shutdown price} is 1 dalen a dozen.

\textbf{4.} $1 + 0.25q = p$, so $q = 4(p - 1)$ for $p > 1$.

\textbf{5.} $q = 6$; revenue 15; cost $8 + 6 + 4.5 = 18.5$; profit $-3.5$ (a loss of 350 dalens a week). It should produce: shut, it would lose its fixed cost, 8.

\textbf{6.} $Q = 160(p - 1)$.

\textbf{7.} $160p - 160 = 560 - 80p$ gives $p = 3$, $Q = 320$, 8 per farm; profit $24 - (8 + 8 + 8) = 0$.

\textbf{8.} No. The cost includes the farmer's time at what she could earn elsewhere, so she earns exactly that; zero \omterm{def:g11:perfect-competition:economic-profit}{economic profit} means she earns no more.

\textbf{9.} $160p - 160 = 680 - 80p$ gives $p = 3.5$, $Q = 400$, 10 per farm.

\textbf{10.} $35 - (8 + 10 + 12.5) = 4.5$, that is 450 dalens a week.

\textbf{11.} Each farm earns more than its owner could earn elsewhere; someone setting up a new farm with the same costs would earn it too.

\textbf{12.} 3 dalens a dozen, the lowest \omterm{def:g11:production-and-costs:atc}{average total cost}.

\textbf{13.} $680 - 240 = 440$.

\textbf{14.} $440/8 = 55$ farms.

\textbf{15.} The price is 3 in both; the quantity rose from 320 to 440; each farm produces 8 in both; the number of farms rose from 40 to 55.

\textbf{16.} The price stays at 3.5, and each farm makes 450 dalens a week. A licence is worth 450 dalens a week to its holder: that is what a buyer would pay, each
week, to hold one.

\textbf{17.} $C = 8 + 1.5q + 0.125q^2$: \omterm{def:g11:production-and-costs:atc}{average total cost} is still lowest at $q = 8$, now at $1 + 1.5 + 1 = 3.5$. The long-run price is 3.5; $Q = 680 - 280 = 400$, so
$400/8 = 50$ farms.

\textbf{18.} Her accounts leave out 3 hundred dalens of cost, so her \omterm{def:g11:perfect-competition:economic-profit}{accounting profit} is 300 dalens a week while her \omterm{def:g11:perfect-competition:economic-profit}{economic profit} is 0.

\textbf{19.} More than 3. A higher rent raises every farm's \omterm{def:g11:production-and-costs:atc}{average total cost}, so the lowest average cost, and with it the long-run price, rises with the number of
farms.

\textbf{20.} After the rise in demand, Brenland's egg industry settles at 55 farms, selling a dozen eggs at 3 dalens, the same long-run price as before.
\end{solution}
